% STATUS: written 2026-11-05 (Suh 1969 pp. 255--258 read in full; Nowicki 1983
% pp. 281--297 read at the introduction, Definition 2.2, Theorem 3.7 and
% Corollary 4.4; Coelho--Milies 1993 pp. 263--267 read in full; Jondrup 1995
% pp. 205--218 read at the introduction and Theorem 2; Alizadeh 2009 pp. 574--578
% read in full.)
% STATUS: extended 2026-11-06 (Ghosh and Prakash 2018 read in full)
% STATUS: extended 2026-11-07 (Xu and Zhang 2010, Huang--Li--Qian 2018,
% Ayupov--Arzikulov 2013 and Ou--Wang--Yao 2007 read in full)
% Sources: Ghosh and Prakash 2018 (pp. 1--5, 8); Xu and Zhang 2010
% (pp. 753--759); Huang, Li and Qian 2018 (pp. 1--4); Ayupov and Arzikulov
% 2013 (pp. 1--2, 8); Ou, Wang and Yao 2007 (pp. 378--381).
\section{Full matrix algebras}\label{sec:matrix-triangular:matrix}

The chapter opens where the literature of the chapter opens, with an algebra that is built
from matrices and whose structure was determined by elementary entry-by-entry
work \cite[p.~255]{Suh1969}. Suh takes a finite-dimensional Lie algebra $L$ over a field of
characteristic not $2$, with an associative symmetric bilinear form $(a,b)$ and
the property $[a[bc]]=(a,c)b-(a,b)c$ \cite[p.~255]{Suh1969}. The set of $2\times2$
vector matrices
\[
\begin{pmatrix} \alpha & a \\ b & \beta\end{pmatrix},
\qquad \alpha,\beta\in\Phi,\ a,b\in L,
\]
with the multiplication displayed on the first page of the paper, is a flexible
algebra over $\Phi$ \cite[p.~255]{Suh1969}. Suh asks which alternative algebras
arise this way and answers completely (Theorem~\ref{thm:matrix:suh}) \cite[p.~255]{Suh1969}.

\begin{theorem}[{\cite[Theorem~1, p.~257]{Suh1969}}]\label{thm:matrix:suh}
Let $L$ be a finite dimensional Lie algebra over a field $\Phi$ of
characteristic not $2$, equipped with an associative symmetric bilinear form
$(a,b)$ and having the property $[a[bc]]=(a,c)b-(a,b)c$. If $(a,b)$ is
nondegenerate, then the algebra $\mathfrak{A}$ built from $L$ is the split
Cayley algebra or a quaternion algebra.
\end{theorem}

Suh adds that a nondegenerate form forces $L$ to be simple of dimension three
or one \cite[Corollary, p.~257]{Suh1969}. For a form that is degenerate the
algebra splits as the direct sum of its radical and a subalgebra that is
quaternion or two-dimensional separable \cite[p.~255]{Suh1969}. The lesson
carries over to the derivation question: a matrix shape turns one problem about
a large algebra into several problems about its entries \cite[p.~255]{Suh1969}.

The same entry-by-entry method settles derivations of matrix rings
\cite[p.~263]{CoelhoMilies1993}. Take a ring
$R$ with identity and the ring $T_n(R)$ of upper triangular $n\times n$ matrices
over $R$ \cite[p.~263]{CoelhoMilies1993}. Coelho and Milies describe every
derivation of $T_n(R)$ under that one hypothesis
\cite[p.~263]{CoelhoMilies1993}.

\begin{theorem}[{\cite[p.~263]{CoelhoMilies1993}}]\label{thm:matrix:coelho}
Let $R$ be a ring with unity and let $d\colon T_n(R)\to T_n(R)$ be a
derivation. Then there exists a derivation $\delta\colon R\to R$ and a matrix
$A\in T_n(R)$ such that $d=\overline{\delta}+d_A$.
\end{theorem}

Here $\overline{\delta}$ sends $(x_{ij})$ to $(\delta(x_{ij}))$, and $d_A$ is
the inner derivation $X\mapsto AX-XA$ \cite[p.~263]{CoelhoMilies1993}. The
authors compute $A$ from the values of $d$ on the matrix units and then recover
$\delta$ from $d$, so the description is effective (Proposition~\ref{prop:matrix:coelho-unique})
\cite[p.~264]{CoelhoMilies1993}.\begin{proposition}[{\cite[p.~266]{CoelhoMilies1993}}]\label{prop:matrix:coelho-unique}
Let $d\colon T_n(R)\to T_n(R)$ be a derivation, and let $A,\delta$ be as in
Theorem~\ref{thm:matrix:coelho}. If $A'\in T_n(R)$ and
$\delta'\colon R\to R$ are such that $d=d_{A'}+\overline{\delta'}$, then there
exists an element $a\in R$ such that $\delta'=\delta-\delta_a$ and
$A'=A+aI$.
\end{proposition}

Jøndrup re-proves the same description for $M_n(A)$ and $T_n(A)$ with a short
induction, and adds the corresponding automorphism theorems (Theorem~\ref{thm:matrix:jondrup})
\cite[p.~205]{Jondrup1995}.\begin{theorem}[{\cite[Theorem~2, p.~209]{Jondrup1995}}]\label{thm:matrix:jondrup}
Let $A$ be a $C$-algebra. Every $C$-derivation of $\operatorname{Mat}_n(A)$ or
$T_n(A)$ is the sum of an inner derivation and a derivation induced from a
$C$-derivation of $A$.
\end{theorem}

The proof reduces to the case where the derivation vanishes on $e_{11}$ and on
every element of $f\operatorname{Mat}_n(A)f$, and shows that such a derivation is
inner \cite[p.~209]{Jondrup1995}. The reduction uses only the idempotents
$e_{ii}$, so no hypothesis on $A$ beyond the $C$-algebra structure is needed.

Nowicki asks a finer question \cite[p.~281]{Nowicki1983}. Fix a reflexive and transitive relation $\rho$ on
$\{1,\dots,n\}$ and let $P$ be the set of matrices in $M_n(R)$ whose $(i,j)$
entry vanishes whenever $(i,j)\notin\rho$ \cite[p.~281]{Nowicki1983}. This $P$
is a subring of $M_n(R)$, called a special subring
\cite[p.~281]{Nowicki1983}. An $R$-derivation is a derivation of $P$ that
vanishes on the scalar matrices \cite[p.~284]{Nowicki1983}. Nowicki describes
all of them (Theorem~\ref{thm:matrix:nowicki}) \cite[p.~287]{Nowicki1983}.

\begin{theorem}[{\cite[Theorem~3.7, p.~287]{Nowicki1983}}]\label{thm:matrix:nowicki}
Let $P$ be a special subring of $M_n(R)$ with the relation $\rho$. Every
derivation $d$ of $P$ has a unique representation
$d=[A,-]+\Delta_f+\Theta_\delta$, where $A$ is a matrix in $P$ with
$A_{pp}=0$ for $p=1,\dots,n$, $\delta=\{\delta_t\colon t\in T_\rho\}$ is a set of
derivations of $R$, and $f\colon\rho\to R$ is a quasi-trivial transitive
mapping.
\end{theorem}

The three pieces have different jobs \cite[p.~287]{Nowicki1983}. The commutator
$[A,-]$ is inner \cite[p.~287]{Nowicki1983}. The term $\Delta_f$ moves each entry
by the value of the transitive mapping $f$ \cite[p.~287]{Nowicki1983}. The term
$\Theta_\delta$ applies a derivation of $R$ to the entries, one derivation for
each equivalence class of $\rho$ \cite[p.~287]{Nowicki1983}. The question of
when every derivation is inner then becomes a question about $\rho$ (Corollary~\ref{cor:matrix:nowicki-inner})
\cite[p.~289]{Nowicki1983}.\begin{corollary}[{\cite[Corollary~4.4, p.~289]{Nowicki1983}}]\label{cor:matrix:nowicki-inner}
The following conditions are equivalent: every $R$-derivation of $P$ is inner,
and the relation $\rho$ is regular over $Z(R)$.
\end{corollary}

Regularity of $\rho$ over an abelian group $G$ means that every transitive
mapping from $\rho$ to $G$ is trivial \cite[Definition~2.2, p.~283]{Nowicki1983}.
For a relation that is not regular the term $\Delta_f$ survives, and a proper
derivation exists \cite[p.~289]{Nowicki1983}. The special subrings therefore
separate the inner derivations from the outer ones by a property of the index
relation alone \cite[p.~289]{Nowicki1983}.

The last question for full matrix algebras is the Jordan one
\cite[p.~574]{Alizadeh2009}. Jacobson and
Rickart showed that a full matrix ring over a $2$-torsion free unital ring has
no proper Jordan derivation, through their theorems on Jordan homomorphisms
\cite[p.~574]{Alizadeh2009}. Alizadeh gives a direct proof for matrices over a
ring, with values in a bimodule (Theorem~\ref{thm:matrix:alizadeh}).

\begin{theorem}[{\cite[Theorem~3.1, p.~575]{Alizadeh2009}}]\label{thm:matrix:alizadeh}
Let $A$ be a unital associative ring and $M$ a $2$-torsion free $A$-bimodule.
Every Jordan derivation from $M_n(A)$ into $M_n(M)$ is a derivation.
\end{theorem}

The proof writes $D$ as a family of maps $D_{ijkl}(a)=[D(a\otimes E_{ij})]_{kl}$
from $A$ to $M$ and checks the Jordan law on each family
\cite[p.~575]{Alizadeh2009}. Alizadeh uses no Jordan homomorphism, which is why his proof
extends the classical result to values in a module \cite[p.~575]{Alizadeh2009}.

The torsion hypothesis can be dropped for the full matrix algebra over a
commutative ring \cite[p.~2]{GhoshPrakash2018}. Ghosh and Prakash prove the
Jordan statement without it, and add the inner form of the derivations (Theorem~\ref{thm:matrix:ghosh-jordan})
\cite[Theorem~3.2, p.~8]{GhoshPrakash2018}.\begin{theorem}[{\cite[Theorem~3.1, p.~5]{GhoshPrakash2018}}]\label{thm:matrix:ghosh-jordan}
Let $C$ be a commutative ring with unity and $M_n(C)$ the full matrix algebra
over $C$, with $n\geq2$. Then every Jordan derivation of $M_n(C)$ into itself
is a derivation.
\end{theorem}

The same two authors then turn from the Jordan law to the derivations
themselves (Theorem~\ref{thm:matrix:ghosh-inner}) \cite[p.~5]{GhoshPrakash2018}.

\begin{theorem}[{\cite[Theorem~3.2, p.~8]{GhoshPrakash2018}}]\label{thm:matrix:ghosh-inner}
Let $C$ be a commutative ring with unity and $M_n(C)$ the algebra of all
$n\times n$ matrices over $C$, with $n\geq2$. Then every derivation of $M_n(C)$
is an inner derivation.
\end{theorem}

The upper triangular matrix algebra over the same ring behaves in the same way
\cite[Theorem~2.1 and Theorem~2.2, pp.~3 and~5]{GhoshPrakash2018}. The proof of
the Jordan statement writes the image of a matrix unit in the matrix units and
compares coefficients \cite[pp.~5--8]{GhoshPrakash2018}.

A ladder matrix keeps the block shape of an upper triangular matrix and opens
some of the blocks on the diagonal \cite[p.~1]{GhimireHuang2015}. Fix a
field of characteristic not $2$ and a ladder $L$ whose blocks dominate the
diagonal \cite[Definition~2.2, p.~2]{GhimireHuang2015}. The matrices that
live on $L$ form a Lie algebra under the commutator
\cite[p.~1]{GhimireHuang2015}. Ghimire and Huang describe its derivations
as a direct sum of two ideals \cite[Theorem~3.1, p.~7]{GhimireHuang2015}.
The first ideal consists of the adjoint action of the block upper triangular
matrices \cite[p.~7]{GhimireHuang2015}. The second consists of the linear maps from $M_L/[M_L,M_L]$ to the
centre \cite[p.~7]{GhimireHuang2015}. In the simplest case every
derivation is an adjoint action plus a sum of traces of the diagonal blocks
\cite[p.~7]{GhimireHuang2015}. When the characteristic is not $2$ or $3$
and the ladder is strongly dominant, every derivation of the commutator algebra
extends to the whole algebra \cite[Theorem~5.3 and Corollary~5.4, p.~18]{GhimireHuang2015}.
So the entry-by-entry method of this section reaches a family of Lie algebras
as well \cite[p.~1]{GhimireHuang2015}.

The strictly upper triangular matrices form a nilpotent Lie algebra, and its
derivations have a complete description with no hypothesis on the ring
\cite[p.~378]{OuWangYao2007}. Let $R$ be a commutative ring with identity and
let $N(n,R)$ be the Lie algebra of strictly upper triangular $n\times n$
matrices over $R$ \cite[p.~378]{OuWangYao2007}. Ou, Wang and Yao build five
families of derivations: the inner ones, the diagonal ones, the central ones,
and two families of extremal ones \cite[pp.~379--380]{OuWangYao2007}. The
diagonal family uses a diagonal matrix, and the central family uses a vector
of $n-3$ ring elements (Theorem~\ref{thm:matrix:strict-upper}) \cite[pp.~379--380]{OuWangYao2007}.

\begin{theorem}[{\cite[Theorem~3.2, pp.~380--381]{OuWangYao2007}}]\label{thm:matrix:strict-upper}
Let $R$ be a commutative ring with identity and let $n>2$.
\begin{enumerate}
\item If $n\geq4$, every derivation $\varphi$ of $N(n,R)$ is uniquely
\[
\varphi=\operatorname{ad}x+\eta_y+\mu_c+\rho^{(12)}_{r,s}+\rho^{(n-1,n)}_{p,q},
\]
where $x\in N(n,R)$, $y\in D(n,R)$, $c\in R^{n-3}$ and $r,s,p,q\in R$ satisfy
$2r=2p=0$.
\item If $n=3$, every derivation $\varphi$ of $N(n,R)$ is uniquely
$\varphi=\operatorname{ad}x+\eta_y+\theta^{(12)}_r+\theta^{(23)}_s$ with
$x\in N(n,R)$, $y\in D(n,R)$ and $r,s\in R$.
\end{enumerate}
\end{theorem}

Here $\eta_y$ is the adjoint action of a diagonal matrix
\cite[p.~379]{OuWangYao2007}. The maps $\rho$ and $\theta$ are the extremal
derivations, and they exist only in the dimensions shown
\cite[pp.~379--380]{OuWangYao2007}.

A left derivation changes the Leibniz rule in a different direction. An additive
map $\mu$ from a ring $R$ into a left $R$-module $V$ is a module left derivation
when $\mu(xy)=x\mu(y)+y\mu(x)$ \cite[p.~753]{XuZhang2010Left}. A module Jordan
left derivation satisfies the weaker law $\mu(x^{2})=2x\mu(x)$
\cite[p.~754]{XuZhang2010Left}. Xu and Zhang settle both questions for full
matrix rings (Proposition~\ref{prop:matrix:left-zero}) \cite[pp.~755--757]{XuZhang2010Left}.

\begin{proposition}[{\cite[Proposition~2.1, p.~755]{XuZhang2010Left}}]\label{prop:matrix:left-zero}
For $n\geq2$, a module left derivation $\mu\colon M_n(R)\to M_n(V)$ must be
zero.
\end{proposition}

The Jordan law gives a nonzero answer on the same ring, and the trace controls
it (Theorem~\ref{thm:matrix:jordan-left}).

\begin{theorem}[{\cite[Theorem~2.6, pp.~756--757]{XuZhang2010Left}}]\label{thm:matrix:jordan-left}
For $n\geq2$, a map $\mu\colon M_n(R)\to M_n(V)$ is a module Jordan left
derivation if and only if there are strong module Jordan left derivations
$\mu_{ij}\colon R\to V$ with $\mu(A)=(\mu_{ij}(\operatorname{tr}A))$ for every
$A\in M_n(R)$. In particular, the only such map is zero when $V$ is $2$-torsion
free.
\end{theorem}

A module Jordan left derivation is strong when $\mu(x^{2})=2\mu(x)=0$ for every
$x$ \cite[p.~754]{XuZhang2010Left}. So the trace gives the argument of each
entry map, and the entries only distribute the values
\cite[pp.~756--757]{XuZhang2010Left}.

The question also has a nonlinear form. A 2-local derivation matches its value at
any two points by one derivation, and it need not be additive
\cite[p.~1]{AyupovArzikulov2013}. Ayupov and Arzikulov take inner witnesses on
a matrix ring over a commutative ring (Theorem~\ref{thm:matrix:two-local-inner})
\cite[Theorem~2.1, p.~2]{AyupovArzikulov2013}.\begin{theorem}[{\cite[Theorem~2.1, p.~2]{AyupovArzikulov2013}}]\label{thm:matrix:two-local-inner}
Let $\Re$ be a commutative associative ring with identity and let $M_n(\Re)$ be
the matrix ring of $n\times n$ matrices over $\Re$, with $n>1$. Then every inner
2-local derivation on $M_n(\Re)$ is an inner derivation.
\end{theorem}

The paper also shows that every derivation of an associative ring extends to a
derivation of its matrix ring \cite[p.~1]{AyupovArzikulov2013}.

Full matrix algebras are the case where the two algebras of a triangular algebra
coincide in size and the bimodule fills the whole off-diagonal block
\cite[p.~119]{Cheung2001}. Section~\ref{sec:matrix-triangular:triangular}
keeps the block shape and sets part of it to zero.

Ghimire and Huang compute the derivations of the Lie algebra of dominant upper
triangular ladder matrices.
That family sits between the strictly upper triangular and the full matrix
algebra \cite[Abstract, p.~1]{GhimireHuang2015}.

For a division ring $D$, a generalized canonical Jordan form is shown to exist for triangularizable matrices that are algebraic over $D$, and the inner derivations of $D$ enter the proof~\cite{WlInnerDerivationsOfDivisionRing2018}.
